\documentclass[10pt,a4paper]{article}

%------------------------------------------------------------------------------
%	REQUIRED PACKAGES AND  CONFIGURATIONS
%------------------------------------------------------------------------------

% PACKAGES FOR TITLES
\usepackage{titlesec}
\usepackage{color}

% PACKAGES FOR LANGUAGE AND FONT
\usepackage[utf8]{inputenc}
\usepackage[english]{babel}
\usepackage[T1]{fontenc} % Font encoding
\usepackage{kpfonts}
\tolerance=1
\emergencystretch=\maxdimen
\hyphenpenalty=10000
\hbadness=10000
\vfuzz=30pt

% PACKAGES FOR IMAGES
\usepackage{graphicx}
\graphicspath{{Images/}} % Path for images' folder
\usepackage{eso-pic} % For the background picture on the title page
\usepackage{subfig} % Numbered and caption subfigures using \subfloat
\usepackage{caption} % Coloured captions
\usepackage{transparent}
\usepackage{wrapfig}

% STANDARD MATH PACKAGES
\usepackage{amsmath}
\usepackage{amsthm}
\usepackage{bm}
\usepackage[overload]{empheq}  % For braced-style systems of equations
\usepackage{matlab-prettifier}
\usepackage{gensymb}

% PACKAGES FOR TABLES
\usepackage{tabularx}
\usepackage{longtable} % tables that can span several pages
\usepackage{colortbl}
\usepackage{xltabular}

% PACKAGES FOR ALGORITHMS (PSEUDO-CODE)
\usepackage{algorithm}
\usepackage{algorithmic}

% PACKAGES FOR REFERENCES & BIBLIOGRAPHY
\usepackage{hyperref} % Adds clickable links at references
\hypersetup{
	colorlinks=true,
	linkcolor=bluePoli,
	anchorcolor=black,
	citecolor=bluePoli,
	filecolor=black,
	menucolor=black,
	runcolor=black,
	urlcolor=black,
}
\usepackage{cleveref}
\usepackage[square, numbers, sort&compress]{natbib} % Square brackets, citing references with numbers, citations sorted by appearance in the text and compressed
\bibliographystyle{plain} % You may use a different style adapted to your field
\usepackage{tocbibind}

% PACKAGES FOR THE APPENDIX
\usepackage[page]{appendix}
\renewcommand{\appendixpagename}{\Large\color{bluePoli}Appendice}

% PACKAGES FOR ITEMIZE & ENUMERATES 
\usepackage{enumitem}

% OTHER PACKAGES
\usepackage{amsthm,thmtools,xcolor} % Coloured "Theorem"
\usepackage{comment} % Comment part of code
\usepackage{fancyhdr} % Fancy headers and footers
\usepackage{lipsum} % Insert dummy text
\usepackage{tcolorbox} % Create coloured boxes (e.g. the one for the key-words)
\usepackage{stfloats} % Correct position of the tables
\usepackage{cuted}
\usepackage{multicol}

%-------------------------------------------------------------------------
%	NEW COMMANDS DEFINED
%-------------------------------------------------------------------------

% EXAMPLES OF NEW COMMANDS -> here you see how to define new commands
\newcommand{\bea}{\begin{eqnarray}} % Shortcut for equation arrays
\newcommand{\eea}{\end{eqnarray}}
\newcommand{\e}[1]{\times 10^{#1}}  % Powers of 10 notation
\newcommand{\mathbbm}[1]{\text{\usefont{U}{bbm}{m}{n}#1}} % From mathbbm.sty
\newcommand{\pdev}[2]{\frac{\partial#1}{\partial#2}}
% NB: you can also override some existing commands with the keyword \renewcommand

%----------------------------------------------------------------------------
%	ADD YOUR PACKAGES (be careful of package interaction)
%----------------------------------------------------------------------------


%----------------------------------------------------------------------------
%	ADD YOUR DEFINITIONS AND COMMANDS (be careful of existing commands)
%----------------------------------------------------------------------------
\newcommand{\twofig}[6]{
\begin{figure}[H]
	\begin{minipage}{0.48\linewidth}
		\centering
		\includegraphics[width=\linewidth]{#1}
		\caption{#2}
		\label{fig:#3}
	\end{minipage}\hfill
	\begin{minipage}{0.48\linewidth}
		\centering
		\includegraphics[width=\linewidth]{#4}
		\caption{#5}
		\label{fig:#6}
	\end{minipage}
\end{figure}
}

\newcommand{\twofigII}[7]{
\begin{figure}[H]
	\begin{minipage}{0.48\linewidth}
		\centering
		\includegraphics[width=#7\linewidth]{#1}
		\caption{#2}
		\label{fig:#3}
	\end{minipage}\hfill
	\begin{minipage}{0.48\linewidth}
		\centering
		\includegraphics[width=#7\linewidth]{#4}
		\caption{#5}
		\label{fig:#6}
	\end{minipage}
\end{figure}
}

\newcommand{\cfig}[4]{
\begin{figure}[H]
	\centering
	\includegraphics[width=#4\linewidth]{#1}
	\caption{#2}
	\label{fig:#3}
\end{figure}
}

\newcommand{\rfig}[4]{
\begin{wrapfigure}{r}{#4\linewidth}
	\centering
	\includegraphics[width=\linewidth]{#1}
	\caption{#2}
	\label{fig:#3}
	\vspace*{-13pt}
\end{wrapfigure}
}

\newcommand{\lfig}[4]{
\begin{wrapfigure}{l}{#4\linewidth}
	\centering
	\includegraphics[width=\linewidth]{#1}
	\caption{#2}
	\label{fig:#3}
	\vspace*{-13pt}
\end{wrapfigure}
}


% Do not change Configuration_files/config.tex file unless you really know what you are doing. 
% This file ends the configuration procedures (e.g. customizing commands, definition of new commands)
% Set the geometric layout of the document
\usepackage{geometry}
\geometry{
  top=2.5cm,
  left = 1.5cm,
  right = 1.5cm,
  bottom=1.5cm,
  headheight= 2.2cm,
  headsep= 0cm,
}
\raggedbottom 

% Create color bluePoli (-> manuale grafica coordinata:  https://www.polimi.it/fileadmin/user_upload/il_Politecnico/grafica-coordinata/2015_05_11_46xy_manuale_grafica_coordinata.pdf)
\definecolor{bluePoli}{cmyk}{0.4,0.1,0,0.4}

% Custom theorem environments
\declaretheoremstyle[
  headfont=\color{bluePoli}\normalfont\bfseries,
  bodyfont=\color{black}\normalfont\itshape,
]{colored}

\captionsetup[figure]{labelfont={color=bluePoli}} % Set colour of the captions
\captionsetup[table]{labelfont={color=bluePoli}} % Set colour of the captions
%\captionsetup[algorithm]{labelfont={color=bluePoli}} % Set colour of the captions

\theoremstyle{colored}
\newtheorem{theorem}{Theorem}[section]
\newtheorem{proposition}{Proposition}[section]

% Enhances the features of the standard "table" and "tabular" environments.
\newcommand\T{\rule{0pt}{2.6ex}}
\newcommand\B{\rule[-1.2ex]{0pt}{0pt}}

% Algorithm description
\newcounter{algsubstate}
\renewcommand{\thealgsubstate}{\alph{algsubstate}}
\newenvironment{algsubstates}{
    \setcounter{algsubstate}{0}%
    \renewcommand{\STATE}{%
    \stepcounter{algsubstate}%
    \Statex {\small\thealgsubstate:}\space}
    }{}
    
% Custom theorem environment
\newcolumntype{L}[1]{>{\raggedright\let\newline\\\arraybackslash\hspace{0pt}}m{#1}}
\newcolumntype{C}[1]{>{\centering\let\newline\\\arraybackslash\hspace{0pt}}m{#1}}
\newcolumntype{R}[1]{>{\raggedleft\let\newline\\\arraybackslash\hspace{0pt}}m{#1}}

% Custom itemize environment
\setlist[itemize,1]{label=$\bullet$}
\setlist[itemize,2]{label=$\circ$}
\setlist[itemize,3]{label=$-$}
\setlist{nosep}

% Create command for background pic
\newcommand\BackgroundPic{% Adding background picture
	\put(230,358){
		\parbox[b][\paperheight]{\paperwidth}{%
			\vfill
			\centering
			\transparent{0.4}
			\includegraphics[width=0.5\paperwidth]{raggiera_polimi.eps}%
			\vfill
}}}

% Set indentation
\setlength\parindent{0pt}

% Set abstract
\renewenvironment{abstract}{\section*{\abstractname}}{}

% Custom title commands
\titleformat{\section}
{\color{bluePoli}\normalfont\Large\bfseries}
{\color{bluePoli}\thesection.}{1em}{}
\titlespacing*{\section}
{0pt}{2ex}{1ex}

\titleformat{\subsection}
{\color{bluePoli}\normalfont\large\bfseries}
{\color{bluePoli}\thesubsection.}{1em}{}
\titlespacing*{\subsection}
{0pt}{2ex}{1ex}

% Custom headers and footers
\pagestyle{fancy}
\fancyhf{}
      
\fancyfoot{}
\fancyfoot[C]{\thepage} % page
\renewcommand{\headrulewidth}{0mm} % headrule width
\renewcommand{\footrulewidth}{0mm} % footrule width

\makeatletter
\patchcmd{\headrule}{\hrule}{\color{black}\hrule}{}{} % headrule
\patchcmd{\footrule}{\hrule}{\color{black}\hrule}{}{} % footrule
\makeatother

% -> Create the header
\chead[C]{
\centering
\begin{tcolorbox}[arc=0pt, boxrule=0pt, colback=bluePoli!60, width=\textwidth, colupper=white]
    \textbf{Simulation of a LEO orbiting microsat on Simulink} \hfill \textbf{\author}  
\end{tcolorbox}
}

%-----------------------------------------------------------------------------

% Insert here the info that will be displayed into your Title page 
% -> title of your work
\renewcommand{\title}{Simulation of a LEO orbiting microsat on Simulink}

% -> author name and surname
\newcommand{\authors}{
\begin{tabular}{rll}
    10723712	&	Marcello Pareschi	&	(BSc Aerospace Engineering - Politecnico di Milano)\\
    10836125	&	Daniele Paternoster	&	(BSc Aerospace Engineering - Politecnico di Milano)\\
    10711624	&	Alex Cristian Turcu	&	(BSc Aerospace Engineering - Politecnico di Milano)\\
    10884250	&	Tamim Harun	Or		&	(BCs)
\end{tabular}
}

% -> MSc course
\newcommand{\course}{Space Engineering}

% -> advisor name and surname
\newcommand{\professor}{Franco Bernelli Zazzera}

% IF AND ONLY IF you need to modify the co-supervisors you also have to modify the file Configuration_files/title_page.tex (ONLY where it is marked)
%\newcommand{\firstcoadvisor}{Name Surname} % insert if any otherwise comment
%\newcommand{\secondcoadvisor}{Name Surname} % insert if any otherwise comment

% -> academic year
\newcommand{\YEAR}{2023-2024}

\begin{document}
\pagenumbering{Roman}

%-----------------------------------------------------------------------------
% TITLE PAGE
%-----------------------------------------------------------------------------

% Do not change Configuration_files/TitlePage.tex (Modify it IF AND ONLY IF you need to add or delete the Co-advisors)
% This file creates the Title Page of the document
% DO NOT REMOVE SPACES BETWEEN LINES!

\let\temp\newpage
\let\newpage\relax
\begin{titlepage}

\AddToShipoutPicture*{\BackgroundPic}

\hspace{-0.6cm}\includegraphics[width=0.6\textwidth]{logo_polimi_ing_indinf.eps}

\vspace{-0.2cm}
\Large{\textbf{\color{bluePoli}{\title}}}\\
\hspace*{\fill}

\vspace{-0.2cm}
\fontsize{0.3cm}{0.5cm}\selectfont \bfseries \textsc{\color{bluePoli} MSc in \course}\\
\hspace*{\fill}

\vspace{-0.2cm}
\fontsize{0.3cm}{0.5cm} \selectfont \bfseries Authors:\par
\vspace{0.2cm} \hspace{0.3cm} \textsc{\textbf{\authors}}\\
\hspace*{\fill}

\fontsize{0.3cm}{0.5cm}\selectfont \bfseries Professor: \textsc{\textbf{\professor}}\\
\hspace*{\fill}

% if only ONE co-advisor is present:
%\vspace{-0.4cm}
%\fontsize{0.3cm}{0.5cm}\selectfont \bfseries Co-advisor: %\textsc{\textbf{\firstcoadvisor}}\\
% if more than one co-advisors are present:
%\vspace{-0.4cm}
%\fontsize{0.3cm}{0.5cm}\selectfont \bfseries Co-advisors: \textsc{\textbf{\firstcoadvisor}}\textsc{\textbf{\secondcoadvisor}}\\

\vspace{-0.4cm}
\fontsize{0.3cm}{0.5cm}\selectfont \bfseries Academic year: \textsc{\textbf{\YEAR}}

\small \normalfont

\vspace{11pt}

\centerline{\rule{1.0\textwidth}{0.4pt}}

\vspace{15pt}

\end{titlepage}
\let\newpage\temp

\thispagestyle{plain} % In order to not show the header in the first page

\vspace{15mm}
\begin{abstract}
\addcontentsline{toc}{section}{Abstract}
\vspace*{5mm}

The following report discusses the attitude dynamics and control of a microsatellite in Sun-synchronous low-Earth orbit. The mission of the satellite is to point the Earth, due to the payload requirement. The simulation was carried out in Simulink environment. The rotational dynamics was modeled through the Euler equations, the kinematics was parametrized through two sets of Euler angles (312 - 313). The environment disturbances that was accounted for were the magnetic field interaction and gravity gradient torque. They were considered the most relevant after a general analysis of all the four main disturbances (SRP, air drag torque, magentic torque, gravity gradient torque). The orbital motion was modeled as a restricted two body problem. 

The on-board sensors were a horizon sensor, a magnetometer and a sun sensor. They were all modeled in Simulink, taking as reference real sensors. The actuators installed on the satellite were a magnetorquer and two reaction wheels, modeled referring to real actuators. 

The control logic implemented on-board considers two algorithms. The first deals with the de-tumbling of the spacecraft using the B-dot control, until a certain condition on the value of the derivative of B is satisfied. The second contemplates the slew and tracking manoeuvre togheter. This last phase is performed through an extension of the PD controller for non-linear dynamics. 


\end{abstract}
\clearpage


%-----------------------------------------------------------------------------

%-----------------------------------------------------------------------------
% INDEX
%-----------------------------------------------------------------------------

{\hypersetup{linkcolor=black}\tableofcontents}

\clearpage

%-----------------------------------------------------------------------------

%{\hypersetup{linkcolor=black}\listoftables}


%-----------------------------------------------------------------------------

%{\hypersetup{linkcolor=black}\listoffigures}

%\clearpage

%-----------------------------------------------------------------------------

\section{Symbols}
\label{sec:symbols}
\setcounter{subsection}{2}
\newsavebox\ltmcbox
\renewcommand{\arraystretch}{1.2}
\renewcommand\tabularxcolumn[1]{m{#1}}

\begin{multicols}{2}
	\raggedcolumns

	\subsection{Framework analysis}
	\vspace{-10pt}
	\small
	\setbox\ltmcbox\vbox{
	\makeatletter\col@number\@ne
	\begin{xltabular}{\linewidth}{cc>{\raggedright\arraybackslash}X}
		$\bm{I}$ & $[kg/m^2]$ & inertia matrix \\
		$a$ & $[km]$ & semi-major axis \\
		$e$ & $[-]$ & eccentricty \\
		$i$ & $[deg]$ & inclination \\
		$\omega$ & $[deg]$ & pericentre anomaly \\
		$\Omega$ & $[deg]$ & RAAN \\
		$\theta$ & $[rad]$ & true anomaly \\
		$n$ & $[rad/s]$ & mean anomaly \\
		$\mathcal{P}$ & $[-]$ & perifocal frame \\
		$\mathcal{N}$ & $[-]$ & inertial frame \\
		$\bm{A_{PN}}$ & $[-]$ & rotation matrix from intertial to perifocal frame\\
		$\bm{r_p}$ & $[km]$ & S/C position in $\mathcal{P}$ frame \\
		$\bm{r_p}$ & $[km]$ & S/C position in $\mathcal{N}$ frame 
	\end{xltabular}
	\unskip
	\unpenalty
	\unpenalty}
	\unvbox\ltmcbox

	\subsection{Dynamics}
	\vspace{-10pt}
	\small
	\setbox\ltmcbox\vbox{
	\makeatletter\col@number\@ne
	\begin{xltabular}{\linewidth}{cc>{\raggedright\arraybackslash}X}
		$\mathcal{B}$ & $[-]$ & body frame \\
		$\bm{x_b}$ & $[-]$ & x direction of the $\mathcal{B}$ frame \\
		$\bm{y_b}$ & $[-]$ & y direction of the $\mathcal{B}$ frame \\
		$\bm{z_b}$ & $[-]$ & z direction of the $\mathcal{B}$ frame \\
		$\bm{\omega}$ & $[rad/s]$ & angular velocity in $\mathcal{B}$ frame \\
		$\bm{\dot{\omega}}$ & $[rad/s]$ & angular velocity rate in $\mathcal{B}$ frame \\
		$\bm{M_d}$ & $[Nm]$ & disturbances torque \\
		$\bm{M_c}$ & $[Nm]$ & control torque 
	\end{xltabular}
	\unskip
	\unpenalty
	\unpenalty}
	\unvbox\ltmcbox

	\subsection{Kinematics}
	\vspace{-10pt}
	\small
	\setbox\ltmcbox\vbox{
	\makeatletter\col@number\@ne
	\begin{xltabular}{\linewidth}{cc>{\raggedright\arraybackslash}X}
		$\bm{s}$ & $[rad]$ & set of Euler angles \\
		$\bm{\dot{s}}$ & $[rad/s]$ & set of Euler angle's rates \\
		$312$ & $[-]$ & sequence of axis' rotation \\
		$313$ & $[-]$ & sequence of axis' rotation \\
		$\phi$ & $[rad]$ & first rotation's angle \\
		$\theta$ & $[rad]$ & second rotation's angle \\
		$\psi$ & $[rad]$ & third rotation's angle \\
		$\bm{A_{BN}}$ & $[Nm]$ & attitude matrix \\
		$tol$ & $[-]$ & tolerance value for singularity check 
	\end{xltabular}
	\unskip
	\unpenalty
	\unpenalty}
	\unvbox\ltmcbox

	\subsection{Disturbances analysis}
	\vspace{-10pt}
	\small
	\setbox\ltmcbox\vbox{
	\makeatletter\col@number\@ne
	\begin{xltabular}{\linewidth}{cc>{\raggedright\arraybackslash}X}
		$\bm{M}$ & $[Nm]$ & disturbance torque \\
		$\bm{D}$ & $[Am^{2}]$ & magnetic dipole from coils or parasitic currents \\
		$\bm{B}$ & $[T]$ & magnetic field \\
		$V$      & $[Tm]$ & magnetic field scalar potential \\
		$ECEF$   & [-] & Earth Centered Earth Fixed frame \\
		$r$ & $[m]$ & radial distance from Earth's centre to S/C position (ECEF) \\
		$\theta$ & $[rad]$ & longitude (ECEF) \\
		$\phi$ & $[rad]$ & co-latitude (ECEF) \\
        $P^{n,m} (\cos \theta)$ & $[-]$ & Gauss normalized associated Legendre functions \\
		$g^{n,m} (t)$ & $[nT]$ & Schmidt semi-normalized spherical harmonic coefficients \\
        $h^{n,m} (t)$ & $[nT]$ & Schmidt semi-normalized spherical harmonic coefficients \\
		$F_i$ & $[N]$ & disturbance force acting on i-th panel \\
		$P$ & $[Pa]$ & Sun's pressure of radiation \\
		$A_i$ & $[Pa]$ & area of i-th panel \\
		$\bm{\hat{S}_B}$ & $[-]$ & Sun's direction in $\mathcal{B}$ frame \\
		$\bm{\hat{N}_{B,i}}$ & $[-]$ &  i-th panel's normal direction in $\mathcal{B}$ frame \\
		$\rho_s$ & $[-]$ & coefficient of specular reflection \\
		$\rho_a$ & $[-]$ & coefficient of absorption \\
		$\rho_d$ & $[-]$ & coefficient of diffusion \\
		$\epsilon$ & $[deg]$ & obliquity of Earth \\
		$C_D$ & $[-]$ & coefficient of drag \\
		$\bm{r_i}$ & $[m]$ & position of centre of action of the force \\
		$\bm{v_{B,i}}$ & $[m/s]$ & relative velocity in $\mathcal{B}$ frame of the i-th panel \\
		$\bm{\hat{v}_{B,i}}$ & $[s^{-1}]$ & direction of relative velocity in $\mathcal{B}$ frame of the i-th panel \\
		$\rho$ & $[kg/m^3]$ & denisty of air \\
		$\rho_0$ & $[kg/m^3]$ & reference denisty of air from US Standard Atmosphere model (1976)\\
		$h_0$ & $[km]$ & reference heigh from US Standard Atmosphere model (1976)\\
		$H$ & $[km]$ & scaling height from US Standard Atmosphere model (1976)\\
		$G$ & $[m^3 kg^{-1} s^{-2}]$ & universal gravitational constant \\
		$m_t$ & $[kg]$ & mass of Earth \\
		$c_i$ & $[-]$ & director cosines of the radial direction in the $\mathcal{B}$ frame \\
	\end{xltabular}
	\unskip
	\unpenalty
	\unpenalty}
	\unvbox\ltmcbox

	\subsection{Sensors}
	\vspace{-10pt}
	\small
	\setbox\ltmcbox\vbox{
	\makeatletter\col@number\@ne
	\begin{xltabular}{\linewidth}{cc>{\raggedright\arraybackslash}X}
		$Np$ & $[rad^2s]$ & noise power \\
		$Ts$ & $[s]$ & sampling time \\
		$\sigma^2$ & $[rad^2]$ & variance of the measurment \\
		$N.S.D.$ & $[nT/ \sqrt{Hz}]$ & noise spectral density \\
		$I$ & $[A]$ & intensity of current \\
		$\alpha$ & $[A/W]$ & coefficient of the sun sensor \\
		$S$ & $[m^{2}]$ & sensor' surface area \\
		$W$ & $[W/m^{2}]$ & intensity of incident radiation\\
		$\theta$ & $[rad]$ & angle of incident light\\
	\end{xltabular}
	\unskip
	\unpenalty
	\unpenalty}
	\unvbox\ltmcbox

	\subsection{Attitude determination}
	\vspace{-10pt}
	\small
	\setbox\ltmcbox\vbox{
	\makeatletter\col@number\@ne
	\begin{xltabular}{\linewidth}{cc>{\raggedright\arraybackslash}X}
		$J$ & $[-]$ & cost function \\
		$\alpha_i$ & $[-]$ & weight on i-th sensor for J function \\
		$\bm{s_i}$ & $[-]$ & unit vector measured by sensor in $\mathcal{B}$ frame\\
		$\bm{v_i}$ & $[-]$ & unit vector calcuated by on-board model in $\mathcal{N}$ frame \\
		$\phi$ & $[rad]$ & error angle between true and measured direction without bias \\
		$E\left[\phi\right]$ & $[rad]$ & expected valure of error angle \\
		$\sigma_{\phi}^2$ & $[rad^2]$ & variance of error angle \\
		$\sigma_{\phi}$ & $[rad]$ & standard deviation of error angle \\
		
	\end{xltabular}
	\unskip
	\unpenalty
	\unpenalty}
	\unvbox\ltmcbox

	\subsection{Actuators}
	\vspace{-10pt}
	\small
	\setbox\ltmcbox\vbox{
	\makeatletter\col@number\@ne
	\begin{xltabular}{\linewidth}{cc>{\raggedright\arraybackslash}X}
		$\bm{D}$ & $[Am^{2}]$ & magnetic dipole from magnetorquers \\
		$\bm{M}$ & $[Nm]$ & control torque \\
		$\bm{h}$ & $[kgm^2/s]$ & angular momentum of wheels + satellite \\
		$\bm{h_r}$ & $[kgm^2/s]$ & relative angular momentum \\
		$\bm{\dot{h}_r}$ & $[kgm^2/s^2]$ & relative angular momentum rate \\
		$\bm{\omega_{max}}$ & $[RPM]$ & saturation of reaction wheels \\
		$\bm{\dot{h}_{max}}$ & $[Nm]$ & maximum available torque for reaction wheels \\
		$\bm{\omega_r}$ & $[rad/s]$ & relative angular velocity of reaction wheel \\
		$\bm{\dot{\omega}_r}$ & $[rad/s^2]$ & relative angular velocity rate of reaction wheel \\
		$\sigma_\omega$ & $[RPM]$ & standard deviation on $\omega_r$
	\end{xltabular}
	\unskip
	\unpenalty
	\unpenalty}
	\unvbox\ltmcbox

	\subsection{Control logic}
	\vspace{-10pt}
	\small
	\setbox\ltmcbox\vbox{
	\makeatletter\col@number\@ne
	\begin{xltabular}{\linewidth}{cc>{\raggedright\arraybackslash}X}
		$\bm{B_B}$ & $[T]$ & magnetic field in $\mathcal{B}$ frame\\
		$\bm{\dot{\hat{B}}_B}$ & $[1/s]$ & magnetic field rate unit direction in $\mathcal{B}$ frame \\
		$k_{\omega}$ & $[kg \cdot m^2/s]$ & gain for detumbling phase \\
		$\bm{\dot{A}_{B,N}}$ & $[1/s]$ & rate of attitude matrix \\
		$\bm{B_N}$ & $[T]$ & magnetic field in $\mathcal{N}$ frame \\
		$\bm{\dot{B}_N}$ & $[T/s]$ & magnetic field rate in $\mathcal{N}$ frame \\
		$\mathcal{LVLH}$ & $[-]$ & local vertical local horizontal frame \\
		$\bm{A_{BL}}$ & $[-]$ & error matrix between $\mathcal{B}$ and $\mathcal{LVLH}$ frame \\
		$\xi$ & $[\deg]$ & angle between orbital plane and magnetic equator \\
		$\bm{\epsilon}$ & $[-]$ & attitude error for non-linear control \\
		$\bm{\omega_{BL}}$ & $[rad/s]$ & angular velocity error between $\mathcal{LVLH}$ and $\mathcal{B}$ frames \\
		$k_1$ & $[kg \cdot m^2/s^2]$ & gain on attitude error for non-linear control \\
		$k_2$ & $[kg \cdot m^2/s]$ & gain on angular velocity error for non-linear control \\
		$\bm{\alpha}$ & $[rad]$ & small Euler angles between $\mathcal{B}$ and $\mathcal{LVLH}$ frame \\
		$\bm{\dot{\alpha}}$ & $[rad/s]$ & small Euler angles rate between $\mathcal{B}$ and $\mathcal{LVLH}$ frame \\
		$\bm{K_p}$ & $[kg \cdot m^2/s^2]$ & gain matrix on $\bm{\alpha}$ for linear control \\
		$\bm{K_d}$ & $[kg \cdot m^2/s]$ & gain matrix on $\bm{\dot{\alpha}}$ for linear control \\
		$K_x$ & $[-]$ & yaw inertial coefficient \\
		$K_y$ & $[-]$ & roll inertial coefficient \\
		$\bm{p_{open}}$ & $[1/s]$ & open-loop poles for the dynamic system \\
		$\bm{p}$ & $[1/s]$ & closed-loop poles for the dynamic system
	\end{xltabular}
	\unskip
	\unpenalty
	\unpenalty}
	\unvbox\ltmcbox

	\subsection{Simulation results}
	\vspace{-10pt}
	\small
	\setbox\ltmcbox\vbox{
	\makeatletter\col@number\@ne
	\begin{xltabular}{\linewidth}{cc>{\raggedright\arraybackslash}X}
		$\bm{B_m}$ & $[T]$ & measured magnetic field \\
		$\lVert \boldsymbol{\dot{\hat{{B}}}_m} \rVert$ & $[1/s]$ & norm of rate of change of direction of $\bm{B_m}$ \\
		$\bm{\omega_{estim}}$ & $[rad/s]$ & onboard estimation of angular velocity
	\end{xltabular}
	\unskip
	\unpenalty
	\unpenalty}
	\unvbox\ltmcbox

\end{multicols}
\pagebreak

\setcounter{page}{1}
\pagenumbering{arabic}

%-----------------------------------------------------------------------------

\section{Requirements}
\label{sec:requirements}

\begin{table}[H]

    \centering
    \begin{tabular}{|c|c|c|c|}
    \hline
    & \textbf{Assigned specification} & \textbf{Modifications}  & \textbf{Motivation for modifications} \\
    \hline
    \textbf{Platform} & Microsat & - & - \\
    \hline
    \textbf{Attitude parameters} & Euler angles & - & - \\
    \hline
    \textbf{Sensor} & Earth Horizon & Magnetic, Sun &
    \makecell{Magnetic because of magnetorquers \\
    Sun because no eclipse during orbit} \\
    \hline
    \textbf{Actuator} & 3 magnetorquers & 2 reaction wheels &
    Underactuation of magnetorquers \\
    \hline
    \end{tabular}
    
    \caption{Mandatory requirements for the project}
    \label{table:requirements}
    
\end{table}

%-----------------------------------------------------------------------------

\section{Environment}
\label{sec:environment}

\subsection{Satellite characterization }
\label{subsec:sat_characterization}

\subsection{Orbit characterization}
\label{subsec:orbit_characterization}

\subsection{Disturbances analysis}
\label{subsec:disturbances_analysis}
In order to make a realistic simulation of the rotating motion of the spacecraft, the environment 
distrubances has to be taken into account. The preliminary study of these external torques is 
crucial for a realistic simulation. In the following paragraphs a brief introduction will be 
done for all the main disturbances, then the simulation of the specific satellite and orbit
will be presented, mainly to choose the two most relevant disturbances. This choice is reasonable 
since there are always two predominant effects of disturbance, while the other can be supposed small 
(usually some order of magnitude smaller, but always depends on the specific case) for this reason
all the disturbances are analyzed. 


\subsubsection{Magnetic Disturbance}
\label{subsubsec:dist_mag}



\subsubsection{SRP Disturbance}
\label{subsubsec:dist_SRP}
SRP radiation torque is the disturbance generated by electromagnetic waves that impacts on
the spacecraft panels and generate a force. These forces acting on some of the panels could give rise to 
a net torque around the center of mass of the spacecraft. Only sun radiation will be considered 
in this case, a more deep analysis should consider infrared Earth radiation and reflected Earth radiation. 
In addition, no eclipse condition will be analyzed during all the simulation, a reasonable assumption for the 
sun-synchronous case orbit.

\subsubsection{Drag Disturbance}
\label{subsubsec:dist_drag}



\subsubsection{Gravity Gradient Disturbance}
\label{subsubsec:dist_GG}

%-----------------------------------------------------------------------------

\section{Dynamics}
\label{sec:dynamics}
The equations of the dynamics rotating body motion used throughout the 
simulation are the Euler equations since rigid body motion assumption is made. 
The set of equations are referred to the principal axis frame 
of the satellite. This frame will be also called reference frame $\mathcal{B}$,
it is described by three unit vectors $\left\{ \boldsymbol{x_b, y_b, z_b,} \right\}$, 
that are in the direction of principal inertia axis. 

\begin{equation}
    \boldsymbol{I\dot{\omega}} + \boldsymbol{\omega} \times \boldsymbol{I\omega} = \boldsymbol{M_d} + \boldsymbol{M_c}
\end{equation}

In the above equation the external torque has been divided in to 2 contributions, 
with clear distinction. $M_d$ describes the disturbance torques that act on the 
spacecraft due to environment and presented in the previous section, while $M_c$ 
is referred to the control torque that the actuators are generating to perform 
the tasks required by the control logic. \\
With particular reference to the Simulink model, two configurations of the satellite 
were considered: undeployed configuration (for detumbling phase) and extended configuration 
(for slew and pointing phases). As a consequence, the mass distribution and hence the inertia
matrix are different in terms of numerical values. This fact has been taken into account by implementing
a logic in the dynamic block of Simulink, that switches between the two matrices using a flag
based on the activation of the De-Tumbling control. This istantaneous switch is not completely realistic 
since the extraction of the panels would require some finite time, and in some way could influence the 
real dynamic of the satellite. Anyhow, for the microsat considered, the retracted configuration allows a 
faster detumbling, and also inertia loads and stresses are reduced on the solar panels.



%-----------------------------------------------------------------------------
% BIBLIOGRAFIA
%-----------------------------------------------------------------------------

\clearpage
\bibliography{bibliography}

%-----------------------------------------------------------------------------

\end{document}